% preamble.tex
% UMC 디지털 연결성 분석 보고서 - 학술 보고서 스타일 LaTeX preamble
% 엔진: XeLaTeX (kotex)
% 컴파일: xelatex -> bibtex -> xelatex -> xelatex

\documentclass[10pt,a4paper,twocolumn]{article}

% ─────────────────────────────────────────────
% 한글 / 다국어
% ─────────────────────────────────────────────
\usepackage[hangul]{kotex}
\usepackage{fontspec}

% 영문 본문 (serif) / 제목·캡션 (sans-serif)
% TeX Gyre 계열은 TeX Live 표준 배포에 포함되어 안전한 fallback.
% TeX Live의 TeX Gyre 폰트를 파일명으로 직접 지정 (fontconfig 등록 불필요)
\setmainfont{texgyretermes-regular.otf}[
  BoldFont = texgyretermes-bold.otf,
  ItalicFont = texgyretermes-italic.otf,
  BoldItalicFont = texgyretermes-bolditalic.otf,
  Ligatures=TeX
]
\setsansfont{texgyreheros-regular.otf}[
  BoldFont = texgyreheros-bold.otf,
  ItalicFont = texgyreheros-italic.otf,
  BoldItalicFont = texgyreheros-bolditalic.otf,
  Scale=0.95
]
\setmonofont{texgyrecursor-regular.otf}[
  BoldFont = texgyrecursor-bold.otf,
  ItalicFont = texgyrecursor-italic.otf,
  Scale=0.9
]

% 한글 폰트 설정
% 1순위: macOS 시스템 폰트 (Apple SD Gothic Neo / Nanum Myeongjo)
% 폰트 부재 시 아래 fallback 주석 참고
% Pretendard Variable (사용자 폰트 디렉토리에 설치됨) + AppleSDGothicNeo fallback
\setmainhangulfont{Pretendard Variable}[
  BoldFont       = Pretendard Variable,
  BoldFeatures   = {RawFeature={+wght=700}},
  AutoFakeSlant  = 0.2,
  Script         = Hangul
]
\setsanshangulfont{Pretendard Variable}[
  BoldFont       = Pretendard Variable,
  BoldFeatures   = {RawFeature={+wght=700}},
  Scale          = 0.95,
  Script         = Hangul
]
\setmonohangulfont{Pretendard Variable}[
  BoldFont       = Pretendard Variable,
  BoldFeatures   = {RawFeature={+wght=700}},
  Scale          = 0.9,
  Script         = Hangul
]
% Fallback 1 (Pretendard 설치 시):
%   \setmainhangulfont{Pretendard-Regular}[BoldFont=Pretendard-Bold]
% Fallback 2 (Noto Sans/Serif CJK KR):
%   \setmainhangulfont{Noto Serif CJK KR}[BoldFont={Noto Serif CJK KR Bold}]
% Fallback 3 (NanumMyeongjo / NanumGothic):
%   \setmainhangulfont{NanumMyeongjo}[BoldFont=NanumMyeongjoBold]

% ─────────────────────────────────────────────
% 페이지 레이아웃
% ─────────────────────────────────────────────
\usepackage[
  a4paper,
  top=20mm,
  bottom=22mm,
  left=18mm,
  right=18mm,
  columnsep=6mm,
  headheight=14pt,
  headsep=8mm,
  footskip=10mm
]{geometry}

\usepackage{microtype}
\usepackage{setspace}
\setstretch{1.15}

% ─────────────────────────────────────────────
% 색상 토큰
% ─────────────────────────────────────────────
\usepackage{xcolor}
\definecolor{primary}{HTML}{1A2332}   % Deep Ocean Navy
\definecolor{accent}{HTML}{C97B3F}    % Warm Terracotta
\definecolor{muted}{HTML}{5B6470}
\definecolor{rule}{HTML}{D6D9DE}
\definecolor{bgsoft}{HTML}{F5F2EC}

% ─────────────────────────────────────────────
% 그래픽 / 표
% ─────────────────────────────────────────────
\usepackage{graphicx}
\graphicspath{{media/}{media/media/}}

\usepackage{float}
\usepackage{subcaption}
\usepackage{booktabs}
\usepackage{longtable}
\usepackage{array}
\usepackage{multirow}
\usepackage{tabularx}
\newcolumntype{Y}{>{\raggedright\arraybackslash}X}
\newcolumntype{C}{>{\centering\arraybackslash}X}

% ─────────────────────────────────────────────
% 캡션
% ─────────────────────────────────────────────
\usepackage[
  font=small,
  labelfont={bf,sf,color=primary},
  textfont=small,
  labelsep=period,
  justification=raggedright,
  singlelinecheck=false,
  skip=4pt
]{caption}
\captionsetup[sub]{font=footnotesize,labelfont={bf,sf}}

% ─────────────────────────────────────────────
% 섹션 스타일링 (titlesec)
% ─────────────────────────────────────────────
\usepackage{titlesec}

\titleformat{\section}
  {\sffamily\bfseries\large\color{primary}}
  {\thesection}{0.6em}{}
  [\vspace{2pt}{\color{rule}\titlerule[0.6pt]}]
\titlespacing*{\section}{0pt}{14pt}{6pt}

\titleformat{\subsection}
  {\sffamily\bfseries\normalsize\color{primary}}
  {\thesubsection}{0.5em}{}
\titlespacing*{\subsection}{0pt}{10pt}{4pt}

\titleformat{\subsubsection}
  {\sffamily\bfseries\small\color{accent}}
  {\thesubsubsection}{0.5em}{}
\titlespacing*{\subsubsection}{0pt}{8pt}{3pt}

% ─────────────────────────────────────────────
% 머리말 / 발 (fancyhdr)
% ─────────────────────────────────────────────
\usepackage{fancyhdr}
\pagestyle{fancy}
\fancyhf{}
\renewcommand{\headrulewidth}{0.4pt}
\renewcommand{\footrulewidth}{0pt}
\fancyhead[L]{\sffamily\small\color{muted}서울 25개 자치구 디지털 연결성 분석}
\fancyhead[R]{\sffamily\small\color{muted}\leftmark}
\fancyfoot[C]{\sffamily\small\thepage}

\fancypagestyle{plain}{
  \fancyhf{}
  \fancyfoot[C]{\sffamily\small\thepage}
  \renewcommand{\headrulewidth}{0pt}
}

% ─────────────────────────────────────────────
% 수학 / 단위
% ─────────────────────────────────────────────
\usepackage{amsmath}
\usepackage{amssymb}
\usepackage{mathtools}
\usepackage{siunitx}
\sisetup{
  group-separator={,},
  group-minimum-digits=4,
  detect-all
}

% ─────────────────────────────────────────────
% 인용 / 리스트 / 따옴표
% ─────────────────────────────────────────────
\usepackage[shortlabels]{enumitem}
\setlist{topsep=2pt,itemsep=1pt,parsep=0pt,leftmargin=1.4em}

\usepackage{csquotes}
\usepackage{url}

% 참고문헌: natbib + apalike
% 선택 이유: References가 비정형 텍스트라 BibTeX 변환 시 가벼운 저자-연도 스타일이 적합.
% biblatex(APA)는 정식 .bib + biber 빌드가 필수라 초기 변환 비용이 크다.
\usepackage[round,authoryear]{natbib}
\bibliographystyle{apalike}
\setcitestyle{aysep={,}}

% ─────────────────────────────────────────────
% 하이퍼링크 / PDF 메타데이터
% ─────────────────────────────────────────────
\usepackage[
  unicode=true,
  bookmarks=true,
  bookmarksnumbered=true,
  bookmarksopen=true,
  pdfencoding=auto,
  colorlinks=true,
  linkcolor=primary,
  citecolor=accent,
  urlcolor=accent,
  filecolor=primary,
  breaklinks=true
]{hyperref}

\hypersetup{
  pdftitle={서울 25개 자치구 디지털 연결성(UMC) 분석 보고서},
  pdfauthor={UMC Project Team},
  pdfsubject={Urban Digital Connectivity Analysis},
  pdfkeywords={디지털 연결성, 서울, 자치구, 다층모형, LLM 텍스트 분석}
}

% ─────────────────────────────────────────────
% 박스 환경 (tcolorbox)
% ─────────────────────────────────────────────
\usepackage[most]{tcolorbox}

% 메타이론 박스: 이론적 배경/개념 정의 강조용
\newtcolorbox{metabox}[1][]{
  enhanced,
  colback=bgsoft,
  colframe=primary,
  boxrule=0.6pt,
  arc=2pt,
  left=8pt,right=8pt,top=6pt,bottom=6pt,
  fonttitle=\sffamily\bfseries\color{primary},
  coltitle=primary,
  title=메타이론 노트,
  #1
}

% 핵심 메시지 박스: 장/절 핵심 결론 강조
\newtcolorbox{keymsg}[1][]{
  enhanced,
  colback=white,
  colframe=accent,
  boxrule=1pt,
  arc=2pt,
  leftrule=4pt,
  left=8pt,right=8pt,top=5pt,bottom=5pt,
  fonttitle=\sffamily\bfseries\color{accent},
  coltitle=accent,
  title=핵심 메시지,
  #1
}

% 인용 블록: 짧은 인용/요약용 사이드바
\newtcolorbox{quoteblock}[1][]{
  enhanced,
  colback=bgsoft,
  colframe=bgsoft,
  boxrule=0pt,
  leftrule=2pt,
  arc=0pt,
  borderline west={2pt}{0pt}{accent},
  left=8pt,right=8pt,top=4pt,bottom=4pt,
  fontupper=\itshape\small,
  #1
}

% ─────────────────────────────────────────────
% 그림 / 표 매크로
% ─────────────────────────────────────────────
% 단일 단(column) 폭 그림
\newcommand{\fig}[3]{%
  \begin{figure}[htbp]
    \centering
    \includegraphics[width=\linewidth]{#1}
    \caption{#2}
    \label{#3}
  \end{figure}%
}

% 2단 폭(페이지 폭) 가로지르는 그림
\newcommand{\figfull}[3]{%
  \begin{figure*}[htbp]
    \centering
    \includegraphics[width=\textwidth]{#1}
    \caption{#2}
    \label{#3}
  \end{figure*}%
}

% 표 예시 (실제 사용 시 복사해서 변형)
% \begin{table}[htbp]
%   \centering\small
%   \caption{자치구별 UMC 지수 상위 5}
%   \label{tab:umc-top5}
%   \begin{tabularx}{\linewidth}{lYS[table-format=2.2]}
%     \toprule
%     순위 & 자치구 & {UMC} \\
%     \midrule
%     1 & 강남구 & 87.42 \\
%     \bottomrule
%   \end{tabularx}
% \end{table}

% ─────────────────────────────────────────────
% 기타 유틸
% ─────────────────────────────────────────────
\newcommand{\umc}{\textsc{umc}}
\newcommand{\llm}{\textsc{llm}}
\newcommand{\gu}[1]{#1구}

\renewcommand{\arraystretch}{1.15}
\renewcommand{\thefootnote}{\arabic{footnote}}

\widowpenalty=10000
\clubpenalty=10000
\flushbottom

% ─────────────────────────────────────────────
% 문서 정보 (main.tex에서 override 가능)
% ─────────────────────────────────────────────
\title{\sffamily\bfseries\color{primary}\Large
서울 25개 자치구 디지털 연결성(UMC) 분석 보고서}
\author{\sffamily\normalsize UMC Project Team}
\date{\sffamily\small\today}
